%-------------------------------------------------------------------------------
%	ABSTRACT PAGE
%-------------------------------------------------------------------------------

\addchaptertocentry{Хураангуй}

\begin{center}
{\scshape\Large \univname\par}
{\scshape\large \facname\par}\vspace{0.5cm}
{\huge\textbf{{Хураангуй}} \par}
\bigskip
{\Large{\ttitle} \par}
\bigskip

{\normalsize \shortname \par}
\addressname
\end{center}

\textit{\textbf{Түлхүүр үгс: \keywordnames}}
\bigskip

Санхүүгийн байгууллагын ажилтнууд өгөгдлийн сангаас мэдээлэл авахдаа SQL бичих чадвартай шинжээчээс хамаардаг.
Байгалийн хэлний асуултыг SQL болгон хөрвүүлэх (Text2SQL) системүүд англи хэлний BIRD, Spider зэрэг жишиг дээр
өндөр үр дүн үзүүлж байгаа ч монгол хэлэнд зориулсан систем, үнэлгээний сан, туршилтын нотолгоо бараг байхгүй.
Монгол хэлний онцлог --- залгамал үгзүй, кирилл ба латин галигийн холимог бичлэг, өгөгдлийн санд монголоор
хадгалагдсан утга --- нь англи хэлэнд боловсруулсан аргуудыг шууд хэрэглэхэд саад болдог.

Энэхүү ажлаар банк бус санхүүгийн байгууллагын (ББСБ) PostgreSQL өгөгдлийн сан дээр монгол асуултад хариулах
Text2SQL системийг судалгааны үр дүнд тулгуурлан бүтээж, хяналттай туршилтаар үнэлэх бүрэн дамжлагыг
хэрэгжүүлэв. Систем нь хоёр нээлттэй эхийн загварыг --- GRPO-оор сургасан Arctic-Text2SQL-R1-7B ба DB судлах алхамтай
(agentic) sqrl-9b --- тус бүрийн сургалтын протоколоор ажиллуулж, дараах давхаргуудтай:
(1) монгол/англи тайлбар, синоним, enum утга бүхий semantic model ба нэр томьёоны сан;
(2) BM25 ба BGE-M3 embedding-ийн хосолсон хайлт (кирилл ба галигийн түлхүүртэй утгын хайлт, нэр томьёо, жишээ);
(3) sqlglot-д суурилсан SQL хамгаалалт, read-only гүйцэтгэгч ба мөрийн түвшний хамгаалалт (RLS);
(4) алдааны мессежээр засах, олон хувилбарын санал хураалт, итгэлээр татгалзах механизм;
(5) гүйцэтгэлийн ба литерал шалгах tool, LLM шүүмжлэгчтэй agentic давталт;
(6) rejection sampling, SFT (LoRA), GRPO-д суурилсан fine-tuning-ийн дамжлага.

Үнэлгээнд зориулж 8 домэйн, 500 асуулт (dev 300, test 200) бүхий монгол eval set бэлтгэв. Асуултыг LLM агентууд
ноороглон бичиж, бие даасан LLM хянагчид шалгасан. Асуулт бүр кирилл, латин галиг, англи хувилбартай. Үүнд хариулах
боломжгүй (62) ба хоёрдмол утгатай (63) асуултыг тусгайлан оруулсан. Нэг хэв загвараас үүссэн асуултуудыг нэг
хуваалтад байлгасан тул dev ба test-ийн хооронд ойролцоо давхардал илрээгүй. Бүх gold SQL-ийг 9 өгөгдлийн сан
дээр гүйцэтгэж шалгасан. Туршилтын төлөвлөгөө 50 тохиргоог (A1--G6) «нэг
удаад нэг өөрчлөлт» зарчмаар EX, test-suite нарийвчлал, TrustSQL, McNemar тестээр харьцуулна. Мөн өмнөх ажлын
(text2sql-thesis) 150 кейсийг шилжүүлж, хоёр системийг ижил нөхцөлд харьцуулах орчин бүрдүүлсэн.

Үнэлгээний дамжлагын зөв байдлыг gold SQL-ийг загварын хариу болгож орлуулах (oracle) аргаар шалгахад A6, B5, G1
тохиргоонд, хоёр загвар, хоёр өгөгдлийн сан дээр EX 100\% гарсан. SQL injection-ий 28, монгол ба латин prompt
injection-ий 56 халдлагын бүгдийг систем хаасан. Загварын бодит оноог GPU сервер дээр (NVIDIA L4, A100) гаргана.
\ifnum\resultsavailable=1
Туршилтын үр дүнг Бүлэг~\ref{Chapter3}-т үзүүлэв.
\else
(Серверийн үр дүн хараахан ороогүй.)
\fi
